% ============================================================
\section{Completion and Selection Algebra}
\label{sec:constraint-compiler}
% ============================================================

\begin{quote}\small
\textbf{What this automates.}
Many derivations in discrete calculus reduce a large rule space to a unique stencil by layering constraints: locality, symmetry, conservation, horizon compatibility, and admissibility.
This section formalizes that reduction as an admissible-completion scheme: the calculus produces a stability bound, and admissibility selects either a unique completion or a certified residual family.
\end{quote}

\paragraph{Logical role of admissibility.}
The admissible-completion scheme is a selector, not the source of dynamics.
Bedrock continuation enters earlier through nonretroactive exactification
(Axiom~\ref{ax:nonretroactive-exactification}): nonterminal incidence must
continue by strict extension. The present section starts only after that
continuation has already been realized as a defect, residual correction, or
required stability bound on the chosen operator surface. Admissibility then
chooses a minimum-burden completion among those already compelled realized
continuations.

\subsection{The completion pipeline}

The following procedure systematizes the reduction from a general rule space to a unique (or certified non-unique) selection.

\begin{quote}
\small
\textbf{Procedure (Admissible completion scheme).}
\begin{enumerate}[label=\textbf{(C\arabic*)}, leftmargin=2.5em, itemsep=0.3em]
\item \textbf{Locality.}
Specify that the rule at each site depends only on a neighborhood of radius $r$.
This restricts to a finite-dimensional parameter space of local rules.

\item \textbf{Symmetry.}
Impose equivariance under a local automorphism group $G$ (e.g., permutations, reflections, rotations of the neighborhood).
This quotients the parameter space by $G$-action.

\item \textbf{Conservation invariants.}
Require that certain linear functionals (e.g., total mass, momentum) are preserved.
Additionally, impose polynomial invariants such as $\det = 1$ or $\Tr = D$ for transfer matrices.
Each invariant is an algebraic constraint.

\item \textbf{Horizon compatibility.}
Require that the rule respects horizon transitivity: $\CCPi{h} \CCPi{h'} = \CCPi{h}$ for $h \le h'$.
This ensures the rule is well-defined across coherence levels.

\item \textbf{Admissibility selection.}
Among remaining candidates, apply the minimum-completion principle against a required stability bound (Definition~\ref{def:admissibility} below).

\item \textbf{Selection certificate.}
If exactly one rule survives, output it.
If multiple rules survive, output the residual degrees of freedom as a ``non-uniqueness certificate.''
If no rule survives, the constraint system is inconsistent.
\end{enumerate}
\end{quote}

\paragraph{Certificate gates.}
\label{rem:certificate-gates}
Three pre-checks can shortcut the pipeline (see Box~\ref{box:certificate-checklist} for the checklist):
\begin{itemize}[itemsep=0.2em]
\item \CertA\ \textbf{(Zero residual):} If $\CCBlockPQ{A} = 0$ or
      $\CCBlockQP{A} = 0$, no Schur correction is needed and the observed
      operator is exact, so the elimination step may be omitted.
\item \CertB\ \textbf{(Solver exactness):} If $\CCCoh{k}(D) = 0$, the Green operator exists, providing canonical integration for the solver layer.
\item \CertC\ \textbf{(Holographic saturation):} If $\Ker(R_k) = \Img(D_{k-1})$, boundary data fully determines bulk states up to exact forms.
\end{itemize}
When applicable, these certificates provide structural guarantees that simplify or eliminate steps (C4)--(C5).

\begin{definition}[Minimum completion admissibility]
\label{def:admissibility}\lean{Compiler.Admissible}
Let $D$ be a self-adjoint ``required lower bound'' (or ``defect-to-compensate'') operator on the observable space, derived from defect propagation (Section~\ref{sec:defect-rules}) and/or elimination residuals (Section~\ref{sec:resolvent-defect}) in the problem at hand.
A candidate completion $Q$ is \emph{admissible} if:
\begin{enumerate}[label=(\roman*)]
\item \textbf{Completion constraint:} $D + Q \succeq 0$ (net quadratic form is bounded below / stable), and
\item \textbf{Minimality:} among candidates satisfying (i) and the structural constraints (C1)--(C4), $Q$ is minimal in PSD order (no strictly smaller PSD completion is feasible).
\end{enumerate}
\end{definition}

\paragraph{Completion constraint is essential.}
\label{rem:completion-essential}
PSD-minimality without a completion constraint may fail to have a minimizer (the infimum may not be attained) or may admit degenerate minima.
The completion constraint $D + Q \succeq 0$ ensures the optimization is well-posed.

\paragraph{Partial order and non-uniqueness.}
\label{rem:psd-partial-order}
The PSD order $\preceq$ is a \emph{partial} order: not all pairs of PSD matrices are comparable.
Consequently, minimal elements need not be unique.
When multiple minimal elements exist, the selection scheme records the
\emph{face} of the PSD cone containing all minimal candidates, namely a
``non-uniqueness certificate.''

\smallskip\noindent
\textbf{Canonical tie-breaker (optional):}
If a unique selection is required, choose among minimal candidates by:
\begin{enumerate}[label=(\alph*)]
\item Minimum trace: $\arg\min_{Q \text{ minimal}} \Tr(Q)$.
\item If traces are equal, minimum Frobenius norm: $\arg\min \|Q\|_F$.
\end{enumerate}
These tie-breakers are not part of the core definition; they may be applied when a unique output is needed.

\paragraph{Relation to defect rules.}
The chain and product rules of Section~\ref{sec:defect-rules} describe how defects transform under operations.
The completion scheme uses these transformations to propagate constraints
through compositions, reducing the viable rule space at each step.

\begin{theorem}[Completion contract: defect-to-bound-to-completion schema]
\label{thm:compiler-contract}\lean{Compiler.compiler_contract}
Let a derivation be built from a finite sequence of operations:
\begin{enumerate}[label=(\alph*), itemsep=0.1em]
\item Horizon restriction $\CCRestr{(\cdot)}{h}$.
\item Composition and linear combination of operators.
\item Defect propagation using the chain/product rules (Section~\ref{sec:defect-rules}).
\item Elimination/resummation steps (e.g., Schur/Feshbach, Section~\ref{sec:resolvent-defect}) producing residual corrections on observables.
\end{enumerate}
Let $D = D^*$ be any self-adjoint operator that lower-bounds (in PSD order) the net instability contributed by the derived defect/residual terms, as formalized by Definition~\ref{def:admissibility}.

Then:
\begin{enumerate}[label=(\roman*), itemsep=0.2em]
\item \textbf{Canonical target:} The completion problem ``find $Q$ in the constrained candidate family such that $D + Q \succeq 0$'' is the canonical stabilization target produced by the derivation.
\item \textbf{Certification:} Any admissible $Q$ (Definition~\ref{def:admissibility}) yields a certified bounded-below net quadratic form on observables.
\item \textbf{Uniqueness or non-uniqueness:} If the admissible set has a unique minimal element, the scheme outputs a unique rule; otherwise it records the minimal face as a non-uniqueness certificate.
\item \textbf{Scalar case:} When the candidate family is a ray $\{\nu L : \nu \ge 0\}$, the unique minimal completion coefficient is given by a Rayleigh quotient characterization.
\end{enumerate}
\end{theorem}

(Proof in Appendix~\ref{sec:appendix-tower}.)

\paragraph{Why this is ``more than the sum''.}
\label{rem:more-than-sum}
The defect grammar makes $D$ compositional (defects propagate through derivations).
Elimination makes $D$ computable on observables (shadow degrees of freedom are integrated out).
Admissibility plus certificates turn the remaining choice into a constrained optimization with explicit outputs (value and tight mode).
The three components interlock to form a single well-posed contract.

\subsection{Asymptotically closing selectors}
\label{subsec:asymptotically-closing-selectors}

\begin{definition}[Minimal scalar repair]\lean{Compiler.MinimalScalarRepair}
\label{def:minimal-scalar-repair}
Let $D$ be a self-adjoint operator on a finite-dimensional Hilbert space
$\mathcal K$. Write $\lambda_{\min}(D)$ for its least eigenvalue and define the
\emph{minimal scalar repair} of $D$ by
\[
\nu(D):=\max\{0,-\lambda_{\min}(D)\}.
\]
Then $D+\nu(D)I_{\mathcal K}\succeq 0$, and $\nu(D)$ is the least nonnegative
scalar with that property.
\end{definition}

\begin{proof}
If $\lambda_{\min}(D)\ge 0$, then $D\succeq 0$ and the least scalar repair is
$0$. If $\lambda_{\min}(D)<0$, then the spectrum of
$D+\nu(D)I_{\mathcal K}$ is obtained by shifting the spectrum of $D$ by
$-\lambda_{\min}(D)$, so the smallest shifted eigenvalue is $0$. Any smaller
nonnegative scalar leaves a negative eigenvalue, hence cannot produce a
positive semidefinite operator.
\end{proof}

\begin{definition}[Operator-norm minimal admissible selector]\lean{Compiler.OperatorNormMinimalSelector}
\label{def:operator-norm-minimal-selector}
Fix a self-adjoint lower-bound operator $D$ on a finite-dimensional Hilbert
space $\mathcal K$, and let
\[
\mathcal Q
\]
be a finite-dimensional linear space of self-adjoint corrections on
$\mathcal K$. A correction
\[
C^\ast\in\mathcal Q
\]
is an \emph{operator-norm minimal admissible selector} if:
\begin{enumerate}[label=(\roman*), leftmargin=2em, itemsep=0.2em]
\item $D+C^\ast\succeq 0$;
\item among all admissible $Q\in\mathcal Q$ satisfying
      $D+Q\succeq 0$, the correction $C^\ast$ minimizes the operator norm
      $\|\cdot\|$.
\end{enumerate}
\end{definition}

\begin{proposition}[Existence of operator-norm minimal selectors]\lean{Compiler.operator_norm_selector_exists}
\label{prop:operator-norm-selector-exists}
With the notation of Definition~\ref{def:operator-norm-minimal-selector},
assume that $\mathcal Q$ contains the identity $I_{\mathcal K}$. Then an
operator-norm minimal admissible selector exists.
\end{proposition}

(Proof in Appendix~\ref{sec:appendix-tower}.)

\begin{theorem}[Minimal admissible selectors vanish in asymptotically closing Schur families]\lean{Compiler.minimal_selectors_vanish}
\label{thm:minimal-selectors-vanish}
Let $(\CCPi{h})_{h\in\bbN}$ be a faithful horizon tower on $\cH$, let
$A:\cH\to\cH$ be bounded, and fix $z\in\mathbb C$ such that the hypotheses of
Theorem~\ref{thm:schur-tower-strong-recovery-general} hold. For each horizon
$h$, let $D_h$ be a self-adjoint lower-bound operator on $\CCPi{h}\cH$, let
\[
\mathcal Q_h
\]
be a finite-dimensional linear space of self-adjoint corrections on
$\CCPi{h}\cH$ containing the identity $I_{\CCPi{h}\cH}$, and let
$C_h^\ast\in\mathcal Q_h$ be an operator-norm minimal admissible selector.
Assume
\[
\nu(D_h)\to 0
\qquad\text{as }h\to\infty.
\]
Then:
\begin{enumerate}[label=(\roman*), leftmargin=2em, itemsep=0.2em]
\item
      \[
      \|C_h^\ast\|\le \nu(D_h)
      \qquad(h\in\bbN);
      \]
\item for every $x\in\cH$,
      \[
      \|C_h^\ast \CCPi{h}x\|\to 0;
      \]
\item the selected effective family
      \[
      \widetilde A_h(z):=\CCAeff{A}{\CCPi{h}}(z)+C_h^\ast
      \]
      satisfies
      \[
      \widetilde A_h(z)x\to Ax
      \qquad\text{for every }x\in\cH.
      \]
\end{enumerate}
\end{theorem}

(Proof in Appendix~\ref{sec:appendix-tower}.)

\subsection{Normal form under symmetry (commutant / isotypic decomposition)}
\label{subsec:normal-form-symmetry}

When the observable subspace $\cH_{\mathrm{obs}} := \CCPi{h} \cH$ carries a symmetry group action, the constraint (C2) can be made completely explicit: $G$-equivariant operators decompose into independent blocks on isotypic components.

\begin{definition}[$G$-equivariant operator and commutant]
\label{def:commutant}\lean{Compiler.IsGEquivariant}\lean{Compiler.IsInCommutant}
Let $G$ be a finite group acting unitarily on $\cH_{\mathrm{obs}}$ via $U: G \to U(\cH_{\mathrm{obs}})$.
An operator $A: \cH_{\mathrm{obs}} \to \cH_{\mathrm{obs}}$ is \emph{$G$-equivariant} if it commutes with the group action:
\[
A \, U(g) = U(g) \, A \quad \text{for all } g \in G.
\]
The set of all $G$-equivariant operators forms an algebra called the \emph{commutant}:
\[
\mathrm{End}_G(\cH_{\mathrm{obs}}) := \{ A : \cH_{\mathrm{obs}} \to \cH_{\mathrm{obs}} \mid A \, U(g) = U(g) \, A \ \forall g \in G \}.
\]
\end{definition}

\begin{theorem}[Normal form via isotypic projectors]
\label{thm:normal-form}\lean{Compiler.normal_form}
Let $P_\rho$ be the isotypic projectors for the irreducible representations $\rho$ of $G$ appearing in $\cH_{\mathrm{obs}}$.
Then any $G$-equivariant operator $A \in \mathrm{End}_G(\cH_{\mathrm{obs}})$ satisfies:
\[
A = \sum_{\rho} P_\rho \, A \, P_\rho.
\]
That is, $A$ is block-diagonal across isotypic components: $P_\rho \, A \, P_{\rho'} = 0$ for distinct irreps $\rho \ne \rho'$.
\end{theorem}

(Proof in Appendix~\ref{sec:appendix-tower}.)

\begin{corollary}[Block structure on isotypic components]
\label{cor:isotypic-block}\lean{Compiler.isotypic_block_structure}
If the $\rho$-isotypic component has the form $\cH_\rho \cong V_\rho \otimes M_\rho$, where $V_\rho$ is the irreducible representation space and $M_\rho$ is the multiplicity space (of dimension $m_\rho$), then on this component:
\[
P_\rho \, A \, P_\rho \cong I_{V_\rho} \otimes B_\rho,
\]
where $B_\rho: M_\rho \to M_\rho$ is an operator on the multiplicity space.
\end{corollary}

\paragraph{Standard result.}
The tensor factorization and the fact that $G$-equivariant maps act as the identity on the irrep factor (Schur's lemma) is a standard result in representation theory.
Over $\bbR$, the endomorphism algebra of an irrep may be $\bbR$, $\mathbb{C}$, or the quaternions $\mathbb{H}$ (Frobenius--Schur indicator); the parameter count adjusts accordingly.

\begin{corollary}[Parameter count]
\label{cor:parameter-count}\lean{Compiler.parameter_count_formula}
In the finite-dimensional case, if $\cH_{\mathrm{obs}}$ decomposes as $\bigoplus_\rho m_\rho \cdot V_\rho$ (where $m_\rho$ is the multiplicity of irrep $\rho$), then:
\[
\dim \mathrm{End}_G(\cH_{\mathrm{obs}}) = \sum_\rho m_\rho^2 \cdot \dim \mathrm{End}_G(V_\rho),
\]
where $\dim \mathrm{End}_G(V_\rho) \in \{1, 2, 4\}$ over $\bbR$ (corresponding to real, complex, or quaternionic type).
\end{corollary}

\begin{definition}[Equivariant coupling space]\lean{Compiler.EquivariantCouplingSpace}
\label{def:equivariant-coupling-space}
Let $G$ act unitarily on finite-dimensional real Hilbert spaces $X$ and $Y$
through representations $U_X$ and $U_Y$. The \emph{$G$-equivariant coupling
space} from $X$ to $Y$ is
\[
\mathrm{Hom}_G(X,Y):=
\{T:X\to Y \mid T\,U_X(g)=U_Y(g)\,T \text{ for all } g\in G\}.
\]
\end{definition}

\begin{theorem}[Finite equivariant coupling decomposition]\lean{Compiler.equivariant_coupling_decomposition}
\label{thm:equivariant-coupling-decomposition}
Let $G$ act unitarily on finite-dimensional real Hilbert spaces $X$ and $Y$.
Write their isotypic decompositions as
\[
X\cong \bigoplus_{\rho} V_{\rho}\otimes M_{\rho}^X,
\qquad
Y\cong \bigoplus_{\rho} V_{\rho}\otimes M_{\rho}^Y,
\]
where the sum runs over the irreducible representations appearing in either
space, and the multiplicity spaces may be zero. Then every
$G$-equivariant coupling $T\in\mathrm{Hom}_G(X,Y)$ decomposes as
\[
T\cong \bigoplus_{\rho} I_{V_{\rho}}\otimes T_{\rho}
\]
with linear maps
\[
T_{\rho}:M_{\rho}^X\to M_{\rho}^Y.
\]
In particular,
\[
\dim \mathrm{Hom}_G(X,Y)
=
\sum_{\rho}
(\dim M_{\rho}^X)(\dim M_{\rho}^Y)\,\dim \mathrm{End}_G(V_{\rho}),
\]
so the equivariant coupling space is finite-dimensional.
\end{theorem}

(Proof in Appendix~\ref{sec:appendix-tower}.)

\begin{corollary}[Shared-channel criterion for equivariant coupling]\lean{Compiler.equivariant_coupling_shared_channel}
\label{cor:equivariant-coupling-shared-channel}
Let $G$ act unitarily on finite-dimensional real Hilbert spaces $X$ and $Y$.
If there is no irreducible representation $\rho$ for which both multiplicity
spaces $M_\rho^X$ and $M_\rho^Y$ are nonzero, then
\[
\mathrm{Hom}_G(X,Y)=0.
\]
Equivalently, a nonzero $G$-equivariant coupling can occur only through a
shared isotypic channel.
\end{corollary}

\begin{proof}
Under the stated hypothesis, every summand in the dimension formula of
Theorem~\ref{thm:equivariant-coupling-decomposition} vanishes. Therefore
\[
\dim \mathrm{Hom}_G(X,Y)=0,
\]
which is equivalent to $\mathrm{Hom}_G(X,Y)=0$.
\end{proof}

\begin{definition}[Symmetry-scalar operator]\lean{Compiler.SymmetryScalarClosure}
\label{def:symmetry-scalar-closure}
Let $A$ be a self-adjoint $G$-equivariant operator on $\cH_{\mathrm{obs}}$, and
write the isotypic decomposition as
\[
\cH_{\mathrm{obs}} \cong \bigoplus_{\rho} V_{\rho}\otimes M_{\rho}.
\]
The operator $A$ is \emph{scalar on the channel $\rho$} if there is a real
number $\lambda_{\rho}$ such that
\[
P_{\rho}AP_{\rho}\cong I_{V_{\rho}}\otimes \lambda_{\rho}I_{M_{\rho}}.
\]
It is \emph{symmetry-scalar} if it is scalar on every isotypic channel.
\end{definition}

\begin{theorem}[Multiplicity-free equivariant operators are symmetry-scalar]\lean{Compiler.multiplicity_free_closures_scalar}
\label{thm:multiplicity-free-closures-scalar}
Let $A$ be a self-adjoint $G$-equivariant operator on $\cH_{\mathrm{obs}}$. If
the $G$-representation on $\cH_{\mathrm{obs}}$ is multiplicity-free, then $A$
is symmetry-scalar.
\end{theorem}

(Proof in Appendix~\ref{sec:appendix-tower}.)

\begin{theorem}[Orthogonal symmetric-tensor scalarization]\lean{Compiler.orthogonal_symmetric_tensor_scalarization}
\label{thm:orthogonal-symmetric-tensor-scalarization}
Let $d\ge 2$, let $\mathrm{Sym}^2(\bbR^d)$ denote the real symmetric
$d\times d$ matrices, and let
\[
L:\mathrm{Sym}^2(\bbR^d)\to \mathrm{Sym}^2(\bbR^d)
\]
be a linear map satisfying
\[
L(QSQ^\top)=Q\,L(S)\,Q^\top
\qquad\text{for every }Q\in O(d),\ S\in \mathrm{Sym}^2(\bbR^d).
\]
Then there are unique scalars $\alpha,\beta\in\bbR$ such that
\[
L(S)
=
\alpha\Bigl(S-\frac{\operatorname{tr}(S)}{d}I\Bigr)
+
\beta\,\frac{\operatorname{tr}(S)}{d}I
\qquad\text{for every }S\in \mathrm{Sym}^2(\bbR^d).
\]
Equivalently, the trace line $\bbR I$ and the traceless carrier
\[
\mathrm{Sym}^2_0(\bbR^d)
:=
\{S\in \mathrm{Sym}^2(\bbR^d)\mid \operatorname{tr}(S)=0\}
\]
are invariant, and $L$ acts by a scalar on each of them.
\end{theorem}

(Proof in Appendix~\ref{sec:appendix-tower}.)

\begin{quote}\small
\textbf{Intuition.}
Locality (C1) restricts the rule to a finite-dimensional parameter space.
Symmetry (C2) further restricts to $G$-equivariant operators.
The normal form turns ``symmetry'' into a concrete low-dimensional parameterization: one free operator $B_\rho$ per isotypic block, acting only on multiplicity coordinates.
\end{quote}

\paragraph{Implementation in the selection algebra.}
In the admissible-completion scheme, step (C2) is implemented by restricting
candidate operators to $\mathrm{End}_G(\cH_{\mathrm{obs}})$ and using this
normal form.
The isotypic projectors $P_\rho$ are computed via the standard character formula.

\begin{theorem}[Symmetry reduction of admissibility and completion]
\label{thm:symmetry-reduction-admissibility}\lean{Compiler.symmetry_reduction}
Let $\cH_{\mathrm{obs}}$ be finite-dimensional with a unitary representation $U: G \to U(\cH_{\mathrm{obs}})$.
Let $D = D^*$ and candidate completions $Q = Q^*$ be $G$-equivariant (i.e., in $\mathrm{End}_G(\cH_{\mathrm{obs}})$).
Let $\{P_\rho\}$ be isotypic projectors with $\cH_{\mathrm{obs}} = \bigoplus_\rho \cH_\rho$ where $\cH_\rho = P_\rho \cH_{\mathrm{obs}}$.

Then:
\begin{enumerate}[label=(\roman*), itemsep=0.2em]
\item \textbf{Block diagonalization:} $D = \bigoplus_\rho D_\rho$ and $Q = \bigoplus_\rho Q_\rho$ where $D_\rho := P_\rho D P_\rho$ and $Q_\rho := P_\rho Q P_\rho$.
\item \textbf{Completion constraint decomposes:}
\[
D + Q \succeq 0 \quad\Leftrightarrow\quad \text{for all } \rho,\ D_\rho + Q_\rho \succeq 0 \text{ on } \cH_\rho.
\]
\item \textbf{PSD minimality is blockwise:} If $Q$ is minimal among feasible candidates in $\mathrm{End}_G(\cH_{\mathrm{obs}})$, then each $Q_\rho$ is minimal among feasible candidates in $\mathrm{End}_G(\cH_\rho)$.
Conversely, if each $Q_\rho$ is minimal in its blockwise feasible set, then $Q = \bigoplus_\rho Q_\rho$ is minimal globally.
\item \textbf{Certificates are blockwise:} A tight mode $x_*$ with $\langle x_*, (D + Q) x_* \rangle = 0$ can be chosen from a single isotypic component.
A minimal face in $\mathrm{End}_G(\cH_{\mathrm{obs}})$ decomposes as a product of minimal faces in each $\mathrm{End}_G(\cH_\rho)$.
\end{enumerate}
\end{theorem}

(Proof in Appendix~\ref{sec:appendix-tower}.)

\paragraph{Computational reduction.}
\label{rem:symmetry-computational-reduction}
Theorem~\ref{thm:symmetry-reduction-admissibility} reduces the admissibility computation to independent subproblems on each isotypic block.
When multiplicities $m_\rho$ are small, each block is low-dimensional and the Rayleigh quotient can be evaluated explicitly.
Certificates for the global problem are assembled from blockwise certificates.

\subsection{Transport-order selection}
\label{subsec:transport-order-selection}

\begin{definition}[Transport-order filtered admissible family]\lean{TransportOrderFilteredFamily}
\label{def:transport-order-filtered-family}
Let $\mathcal F$ be a family of self-adjoint candidate operators on a
band-split observed sector. The family is \emph{transport-order filtered} if
every $A\in\mathcal F$ has finite tower transport order in the sense of
Definition~\ref{def:tower-transport-order}. Write
\[
\mathcal F_{\mathrm{adm}}\subseteq \mathcal F
\]
for the admissible subset selected by the structural constraints.
\end{definition}

\begin{definition}[Distinguished admissible transport-order class]\lean{distinguishedTransportOrderClass}
\label{def:distinguished-transport-order-class}
Assume the setting of Definition~\ref{def:transport-order-filtered-family} and
suppose $\mathcal F_{\mathrm{adm}}\neq\varnothing$. The
\emph{distinguished admissible transport-order class} is the least integer
\[
m_\ast:=\min\{\operatorname{ord}_{\mathcal T}(A):A\in\mathcal F_{\mathrm{adm}}\},
\]
where $\operatorname{ord}_{\mathcal T}(A)$ denotes the tower transport order
of $A$.
\end{definition}

\begin{theorem}[Minimal admissible transport order is constructive and canonical]\lean{minimal_admissible_transport_order}
\label{thm:minimal-admissible-transport-order}
Assume the setting of Definition~\ref{def:transport-order-filtered-family} and
suppose $\mathcal F_{\mathrm{adm}}\neq\varnothing$. Then:
\begin{enumerate}[label=(\roman*), leftmargin=2em, itemsep=0.2em]
\item the distinguished admissible transport-order class $m_\ast$ of
      Definition~\ref{def:distinguished-transport-order-class} exists;
\item the subclass
      \[
      \mathcal F_{\mathrm{adm}}^{\min}
      :=
      \{A\in\mathcal F_{\mathrm{adm}}:\operatorname{ord}_{\mathcal T}(A)=m_\ast\}
      \]
      is canonical;
\item if $\mathcal F_{\mathrm{adm}}^{\min}$ is singleton modulo the chosen
      symmetry or horizon-preserving equivalence relation, then the transport
      order is forced.
\end{enumerate}
\end{theorem}

(Proof in Appendix~\ref{sec:appendix-tower}.)

\begin{corollary}[Minimal admissible transport order controls exact Schur jets]\lean{minimal_transport_order_controls_jet}
\label{cor:minimal-transport-order-controls-jet}
Assume the setting of Theorem~\ref{thm:minimal-admissible-transport-order}, and
suppose each admissible candidate $A\in\mathcal F_{\mathrm{adm}}$ admits an
exact Schur transport jet as in Definition~\ref{def:schur-transport-jet}. Then
every candidate in the minimal-order subclass
$\mathcal F_{\mathrm{adm}}^{\min}$ has Schur jet coefficients satisfying
\[
\operatorname{ord}_{\mathcal T}(K_{h,n}(A))\le (n+1)m_\ast
\qquad(n\ge 0)
\]
at every horizon where the jet is defined.
\end{corollary}

\begin{proof}
If $A\in\mathcal F_{\mathrm{adm}}^{\min}$, then
\[
\operatorname{ord}_{\mathcal T}(A)=m_\ast.
\]
The coefficient bound is therefore exactly
Theorem~\ref{thm:schur-transport-jet-order} with $r=m_\ast$.
\end{proof}

\begin{theorem}[Order-one forcing criterion]\lean{order_one_forced_transport_order, order_one_forcing_forced_class}
\label{thm:order-one-forcing-criterion}
Assume the setting of
Theorem~\ref{thm:minimal-admissible-transport-order}. Suppose moreover that:
\begin{enumerate}[label=(\roman*), leftmargin=2em, itemsep=0.2em]
\item there exists an admissible candidate $A_1\in\mathcal F_{\mathrm{adm}}$
      with
      \[
      \operatorname{ord}_{\mathcal T}(A_1)\le 1;
      \]
\item no admissible candidate has transport order $0$.
\end{enumerate}
Then the distinguished admissible transport-order class is
\[
m_\ast=1.
\]
If the minimal-order admissible subclass is singleton modulo the chosen
equivalence relation, then the one-step transport class is forced.
\end{theorem}

\begin{proof}
By hypothesis (i),
\[
m_\ast \le \operatorname{ord}_{\mathcal T}(A_1)\le 1.
\]
By hypothesis (ii), one cannot have $m_\ast=0$. Therefore $m_\ast=1$. The final
statement is then exactly part (iii) of
Theorem~\ref{thm:minimal-admissible-transport-order}.
\end{proof}

\begin{definition}[Transport-generated admissible family]\lean{TransportGeneratedAdmissibleFamily}
\label{def:transport-generated-admissible-family}
Assume the setting of Definition~\ref{def:transport-order-filtered-family}. The
admissible family is \emph{transport-generated} if:
\begin{enumerate}[label=(\roman*), leftmargin=2em, itemsep=0.2em]
\item the hypotheses of Theorem~\ref{thm:order-one-forcing-criterion} hold;
\item the minimal-order admissible subclass
      \[
      \mathcal F_{\mathrm{adm}}^{\min}
      \]
      is singleton modulo the chosen equivalence relation and is represented by
      a bounded transport base
      \[
      B:F\to F;
      \]
\item the representative $B$ is without preferred direction at quadratic order
      in the sense of
      Definition~\ref{def:no-preferred-direction-quadratic-transport}.
\end{enumerate}
\end{definition}

\begin{theorem}[Transport generation forces a scalar recursion channel]\lean{transport_generation_forces_scalar_recursion}
\label{thm:transport-generation-forces-scalar-recursion}
Assume the setting of
Definition~\ref{def:transport-generated-admissible-family}, and let
\[
B:F\to F
\]
be the bounded representative of the unique minimal-order admissible
equivalence class. Then:
\begin{enumerate}[label=(\roman*), leftmargin=2em, itemsep=0.2em]
\item the distinguished admissible transport-order class is
      \[
      m_\ast=1;
      \]
\item there exists a unique scalar
      \[
      \kappa(B)\ge 0
      \]
      such that
      \[
      B^*B=\kappa(B)\,I_F;
      \]
\item every bi-infinite sequence
      \[
      x=(x_n)_{n\in\bbZ}
      \]
      propagated by $B$ satisfies the canonical recursion law
      \[
      \mathcal R_Bx=0,
      \qquad\text{equivalently}\qquad
      x_{n+1}=Bx_n-x_{n-1}.
      \]
\end{enumerate}
In particular, the unique minimal admissible transport datum determines a
forced scalar quadratic recursion channel.
\end{theorem}

(Proof in Appendix~\ref{sec:appendix-tower}.)

\subsection{Worked example: transfer matrix selection}

The following minimal example illustrates the completion-and-selection algebra in action.

\paragraph{Worked example (Two-branch selection via trace and determinant).}
\label{ex:transfer-compiler}
Consider a $2 \times 2$ local transfer matrix $M$ subject to:
\begin{itemize}[itemsep=0.2em]
\item \textbf{(C3) Conservation:} $\det(M) = 1$ (unit determinant, invertibility).
\item \textbf{(C3) Trace constraint:} $\Tr(M) = D$ for a fixed invariant $D > 2$.
\end{itemize}
The characteristic polynomial is forced:
\[
\lambda^2 - D\lambda + 1 = 0
\quad\Rightarrow\quad
\lambda_\pm = \frac{D \pm \sqrt{D^2 - 4}}{2}.
\]
\textbf{(C5) Admissibility:} The stable branch $\lambda_+$ (with $\lambda_+ > 1$) is selected; the contracting branch $\lambda_-$ corresponds to an ill-posed (unbounded below) completion.

\smallskip\noindent
\textbf{Output:} Unique admissible eigenvalue $\lambda_+$.

\paragraph{Selection algebra as proof assistant.}
In applications, the admissible-completion scheme is mechanical: once the
constraints (C1)--(C4) are verified, the derivation of the unique rule follows
from the resulting completion problem.
The selection algebra converts domain knowledge (symmetries, conservation laws)
into algebraic constraints, then applies admissibility as the final selection
principle.
